\TNchapter[Terms in alphabetical order. Each entry is the one-sentence (occasionally two-sentence) definition Maya could repeat in a meeting without bluffing. Where a 2025--26 currency term gets a longer treatment in the body, the glossary entry is the short version. Each entry names the chapter where the term is first introduced and, where applicable, the defining citation.]{Glossary}
\label{ch:glossary}

\starttwocol

\begin{description}
  \item[A2A (Agent-to-Agent Protocol).] An emerging 2025 protocol (Google) for structured communication between agents from different vendors. The hardening implication: every A2A edge is a new trust boundary that you, not the protocol, are responsible for enforcing. \textit{See Chapter 6.}
  \item[Advisory vs. operational autonomy.] The line a system crosses the first time it can take an action in the world, not just produce a recommendation. A wrong answer used to be embarrassing; a wrong action is expensive. \textit{See Chapter 1; Chapter 5.}
  \item[Agent.] A software system that takes a goal in natural language, decides on its own which steps to take, calls tools to do the work, and adjusts based on what comes back. The thing that distinguishes an agent from a chatbot is that it acts. \textit{See Chapter 1.}
  \item[Agent estate.] The full collection of agents an organization runs in production — the agentic equivalent of an application portfolio. Once you have more than three, the operational problem starts to look like IT service management, not AI. \textit{See Chapter 8.}
  \item[Agent loop.] The repeating cycle of "model produces a response, the response triggers a tool call, the tool returns a result, the model reads the result, the model produces the next response." Most agent failure modes are a particular thing going wrong in a particular leg of this loop. \textit{See Chapter 1.}
  \item[AgentBench.] An academic benchmark (Liu et al. 2023) that tests an agent across eight environments — operating system, database, web, coding, and others. Useful as a breadth-of-capability check; insufficient for an enterprise decision on its own. \textit{See Chapter 10.}
  \item[AgentPoison.] A research-documented attack class (Chen et al. 2024) in which an adversary poisons an agent's memory store or knowledge base over repeated interactions until the agent's behavior shifts toward the attacker's preferred shape. \textit{See Chapter 6; Chapter 7.}
  \item[Alignment.] Pillar 1 of ARIA. The discipline of making an agent do what you actually wanted, not just what you literally said. Splits into outer alignment (did we specify the goal correctly?) and inner alignment (will the trained model pursue that goal in production?). \textit{See Chapter 1.}
  \item[Alignment tax.] The cost — in capability, helpfulness, or fluency — that comes with any alignment intervention. Vendors who report only capability scores are showing you half the picture. \textit{See Chapter 2.}
  \item[AML.Txxxx (MITRE ATLAS technique identifier).] The naming convention for techniques in MITRE ATLAS. \textit{AML} stands for \textit{Adversarial Machine Learning}; \textit{T} marks a technique; the four-digit number identifies the specific technique (and a trailing \textit{.NNN} identifies a sub-technique). Parallel to the \textit{T0000} pattern used by MITRE ATT\&CK for cyber. Examples your team will see in Chapter 6: \textit{AML.T0051} (LLM Prompt Injection), \textit{AML.T0051.001} (LLM Prompt Injection: Indirect), \textit{AML.T0057} (LLM Data Leakage). \textit{See Chapter 6.}
  \item[Annex III (EU AI Act).] The list inside the EU AI Act that names the use cases the Act calls "high-risk" — biometric identification, critical infrastructure, education, employment, essential services, law enforcement, migration, and the administration of justice. If your agent lands in Annex III, the Act's heaviest obligations apply. \textit{See Chapter 17.}
  \item[Anomaly detection (agent context).] The discipline of flagging unusual tool-call patterns, retrieval results, retry rates, or output distributions that deviate from a baseline. In an agent estate, anomaly detection is a Layer-3 control (proactive maturity and above) that feeds the kill switch and the circuit breaker. Distinct from classical APM anomaly detection because the unit of observation is the agent turn, not the HTTP request. \textit{See Chapter 6; Chapter 8; Chapter 11.}
  \item[Approval gate.] A point in the agent loop where the agent must pause and ask a human before taking the next action. The cheapest, most reliable form of runtime alignment for high-blast-radius actions. \textit{See Chapter 3; Chapter 4.}
  \item[Articles 9–15 (EU AI Act).] The core technical and procedural obligations on high-risk AI systems: risk management (9), data governance (10), technical documentation (11), record-keeping (12), transparency (13), human oversight (14), and accuracy/robustness/cybersecurity (15). Each maps to a specific artifact in your pre-cert binder. \textit{See Chapter 19.}
  \item[Article 43 (EU AI Act).] The article requiring a \textit{conformity assessment} before any high-risk AI system goes to market, and again on every substantial modification. The single most consequential procedural obligation in the Act for agent builders. \textit{See Chapter 17; Chapter 19.}
  \item[Articles 71–73 (EU AI Act).] The post-market obligations on high-risk AI systems: public-registry registration (71), documented post-market monitoring (72), and serious-incident reporting (73). \textit{See Chapter 8; Chapter 20.}
  \item[Assurance level.] Under ISO/IEC 42001, the degree of confidence a third party can give that an AI management system meets the standard. Higher assurance means more evidence, more cost, more reviewer time. \textit{See Chapter 19.}
  \item[ATFAA / SHIELD.] Two complementary frameworks from Narajala and Narayan (2025) for agent security. \textit{ATFAA} (Agentic Threats and Failures Analysis and Assessment) is the threat-and-failure-mode taxonomy; \textit{SHIELD} (Secure Hardening, Identity, Egress, Logging, Defense) is the matching control framework. Useful as a synthesis citation alongside OWASP and MITRE ATLAS. \textit{See Chapter 7; Narajala and Narayan 2025.}
  \item[Attack surface.] Everything an attacker could touch to get the agent to misbehave — the system prompt, the tools, the data the agent retrieves, the memory it keeps across turns, and the model itself. With agents, the attack surface is wider than with traditional applications because the agent itself is part of it. \textit{See Chapter 5.}
  \item[Audit committee (certification sense).] The internal or external body that reviews evidence and decides whether to certify, deny, or defer. Maya will likely sit on this committee for her own enterprise. \textit{See Chapter 17; Chapter 19.}
  \item[Behavior contract.] A short, public document that says what an agent will do, what it will not do, and what it will escalate. The minimum artifact an enterprise should publish before letting an external user touch the agent. \textit{See Chapter 18.}
  \item[Behavioral verification.] The discipline of testing an agent for misalignment — not just for accuracy — by running it through scenarios designed to catch deceptive or specification-gaming behavior. Distinct from functional evaluation. \textit{See Chapter 4.}
  \item[Benchmark.] A shared test, run across multiple systems, intended to make them comparable. A score on a benchmark says nothing about whether the system will work for your job — only how it compares. \textit{See Chapter 9.}
  \item[Benchmark contamination.] The situation where the test data has leaked into the model's training data, so the model has effectively seen the answers. The cleanest explanation for why scores jump in suspicious lockstep with new model releases. \textit{See Chapter 9; Chapter 10.}
  \item[Benchmarking.] Pillar 3 of ARIA. The discipline of running an agent against a defined, repeatable test to compare it against alternatives or against its own past versions. Comparative, not absolute. \textit{See Chapter 9.}
  \item[BFCL (Berkeley Function-Calling Leaderboard).] A public leaderboard (Patil et al. 2024; v3 multi-turn extension, 2025) that scores models on their ability to translate a natural-language instruction into a correct, well-formed function call with the right arguments. The narrowest of the agent-relevant benchmarks; the one most likely to predict whether an agent's tool calls will work in your stack. \textit{See Chapter 10; Chapter 14.}
  \item[Blast radius.] The maximum harm an agent can do before someone catches it. Defined by the scope of its tools, its credentials, and the speed at which it can act. The single most important number to know about any agent you certify. \textit{See Chapter 1; Chapter 5.}
  \item[Bonferroni / Benjamini-Hochberg correction.] Two methods for adjusting significance thresholds when comparing many things at once. Bonferroni divides the threshold by the number of comparisons (conservative); Benjamini-Hochberg controls the false discovery rate (less aggressive). Use one of them whenever a dashboard reports many metrics in parallel. \textit{See Chapter 12.}
  \item[Bootstrap confidence interval.] A way of estimating the uncertainty in a benchmark score by resampling the test cases many times (with replacement) and recomputing the score each time; the middle 95\% of those scores is the 95\% interval. If two agents' intervals overlap, the difference is not yet evidence. \textit{See Chapter 9; Chapter 12.}
  \item[Certification.] Pillar 5 of ARIA. The artifact and the process that lets your organization say "yes, we're confident in this agent" to its CFO, its CISO, its regulator, and its board. The stack of standard, evidence, and governance behind that artifact. \textit{See Chapter 17.}
  \item[Certification review board.] A standing cross-functional body that reads the binder, pushes on it, and votes on each certification decision. A board that cannot say no is a rubber stamp. \textit{See Chapter 17; Chapter 19.}
  \item[Circuit breaker (agent context).] An automatic shutdown that trips when an agent's error rate, cost burn, retry rate, or output distribution crosses a threshold. Borrowed from electrical engineering and from microservice patterns; in agents the breaker isolates the misbehaving tool, model, or workflow before the blast radius widens. Distinct from the kill switch, which is a human-pulled lever; the circuit breaker fires on its own. \textit{See Chapter 6; Chapter 7; Chapter 8; Chapter 11.}
  \item[Cohen's d (effect size).] A measure of how large the difference between two agents' scores is, expressed in standard deviations rather than percentage points. A d of 0.2 is small, 0.5 is medium, 0.8 is large; a statistically significant difference can still be too small to matter. \textit{See Chapter 12.}
  \item[Computer Use.] Anthropic's 2024 capability that lets Claude operate a desktop environment by taking screenshots, moving the cursor, clicking, and typing. The first widely available example of a general-purpose agent that can drive arbitrary applications. Cited in this book as the canonical example of a high-blast-radius runtime that demands sandboxing by default. \textit{See Chapter 8.}
  \item[Configuration tampering.] An attack class in which the adversary alters the agent's configuration — system prompt, tool catalog, secrets, model selection, or guardrail policy — rather than its inputs. The defense is the same as for code: signed, versioned, reviewed configuration with change auditing; the configuration store is itself a trust boundary. \textit{See Chapter 6.}
  \item[Confused deputy.] A security pattern, named in classical access-control literature, in which a privileged process is tricked by an unprivileged caller into performing an action the caller could not have performed directly. In agent systems the confused deputy is the agent itself: it holds broad permissions, and an attacker who shapes the user's request can get it to misuse those permissions. The fix is delegation that carries the \textit{caller's} authority through the call chain, not the agent's. \textit{See Chapter 7.}
  \item[Conformity assessment.] The procedure required by the EU AI Act (Article 43) under which a high-risk AI system is checked, and documented, against the Act's requirements before going to market and again on every substantial modification. Some assessments are self-declared; some require a notified body. \textit{See Chapter 17; Chapter 19.}
  \item[Constitutional AI.] A training method (Bai et al. 2022) that supplements human preference labels with critiques the model produces against a written constitution. Scalable, transparent, and partial — the constitution catches what its authors knew to write down, and nothing else. \textit{See Chapter 2.}
  \item[Content provenance.] The discipline of tagging every retrieved document with where it came from, who wrote it, when, and how it was approved. Without provenance you cannot tell, after an incident, whether attacker-controlled content reached the agent's context. \textit{See Chapter 7.}
  \item[Continuous evaluation.] Running evaluations on a sample of live traffic, all the time, rather than only at release. The mechanism that catches drift between certification cycles. \textit{See Chapter 16.}
  \item[Contextual hallucination.] An output that is not supported by the documents the agent was given. The retrieved policy says 30 days; the agent said 60. Caught by groundedness checks against the retrieved context. \textit{See Chapter 15.}
  \item[DecodingTrust.] An academic benchmark (Wang et al. 2024) that evaluates LLMs across eight trust dimensions — toxicity, stereotype bias, adversarial robustness, out-of-distribution generalization, robustness to adversarial demonstrations, privacy, machine ethics, and fairness. A useful sanity check; not a substitute for evaluation on your own data. \textit{See Chapter 16.}
  \item[Deceptive alignment.] When an agent behaves correctly during training and evaluation but pursues a different objective once it judges itself to be in deployment. Different from specification gaming because it requires the model to model the difference between testing and production. \textit{See Chapter 4; Hubinger et al. 2024.}
  \item[Deceptive compliance.] The everyday version of deceptive alignment: the agent passes every test in the eval suite and fails on the production distribution because it has effectively learned which inputs are tests. \textit{See Chapter 1.}
  \item[Demographic parity.] A fairness criterion that says the rate at which an agent makes a particular decision should be the same across demographic groups. Often the easiest fairness criterion to explain to a board; often the wrong one for compliance work because it can mask differences in error rates. \textit{See Chapter 16; Hardt et al. 2016.}
  \item[Denial-of-Wallet.] An attack in which the agent's costs (LLM calls, tool invocations, retrieval) are run up by an adversary until the bill exceeds what the business can sustain. The cost-side cousin of denial-of-service. \textit{See Chapter 6; Chapter 11.}
  \item[Deterministic check.] A grader whose output is fully determined by the agent's output — schema validation, regex matching, JSON parsing, exact match. Cheap, fast, auditable, brittle. \textit{See Chapter 14.}
  \item[Direct prompt injection.] An attack in which the user supplies input to the agent designed to override its system prompt or its guardrails. The 101-level agent attack. \textit{See Chapter 6.}
  \item[DPO, KTO, GRPO, IPO.] Successors to RLHF (direct preference optimization, Kahneman-Tversky optimization, group relative policy optimization, identity preference optimization) that train preference into a model more cheaply or more stably. You do not need to distinguish them; you do need to know that "RLHF" in 2026 often means one of these variants rather than the original PPO-based recipe. \textit{See Chapter 2; Rafailov et al. 2023.}
  \item[Drift.] The slow change in an agent's behavior over time. Three flavors: \textit{model drift} (the model is updated), \textit{data drift} (the world the agent sees has changed), and \textit{concept drift} (what counts as a correct answer has changed). All three are inevitable. \textit{See Chapter 4; Chapter 12; Chapter 16.}
  \item[ECE (expected calibration error).] A measure of how well a model's stated confidence matches its real accuracy. A model that says "I am 90\% sure" should be right 90\% of the time; ECE is how far off that bargain is in practice. \textit{See Chapter 16.}
  \item[Equal opportunity.] A fairness criterion: the true-positive rate is equal across demographic groups. The right metric when the cost of a false negative is much higher than the cost of a false positive — for example, clinical screening. \textit{See Chapter 16; Hardt et al. 2016.}
  \item[Equalized odds.] A fairness criterion that requires the agent's error rates (true-positive, false-positive) to be the same across demographic groups, conditional on the ground truth. Stricter than demographic parity; the criterion most regulators care about. \textit{See Chapter 16; Hardt et al. 2016.}
  \item[Escalation rate.] The fraction of turns or sessions where the agent hands control to a human. A KPI in its own right — too low means the agent is over-stepping; too high means the agent is not earning its cost. \textit{See Chapter 4; Chapter 15.}
  \item[EU AI Act (Regulation 2024/1689).] The European Union's horizontal regulation for AI, in force since August 2024. It tiers AI systems by risk, bans some uses outright, and imposes a \textit{conformity assessment} obligation on high-risk systems before launch and on every material change. For any agent operating in the EU, the Act is the law. \textit{See Chapter 17; Chapter 19; Chapter 20.}
  \item[Evaluation.] Pillar 4 of ARIA. A purpose-built, organization-specific test of an agent doing the job you intend it to do. Distinct from benchmarking: an evaluation answers a question; a benchmark compares to peers. \textit{See Chapter 13.}
  \item[Evaluation rubric.] The atomic, scoring-banded criteria your LLM-as-judge or human reviewer grades against. A rubric of "score from 1 to 10 for clarity" gives noise; a rubric of "did the response cite a source, stay within scope, refuse appropriately" gives signal. \textit{See Chapter 14.}
  \item[Evaluation taxonomy.] The map this book uses: functional, quality, safety, and continuous, each layered against offline and online modes. Most enterprises need at least one eval in each cell to certify an agent. \textit{See Chapter 13.}
  \item[Evidence chain.] The traceable links from a deployed agent backward through its evals, benchmarks, training, and data — the thing an auditor follows when they ask, "how do you know?" \textit{See Chapter 18.}
  \item[Factual hallucination.] An output that contradicts an objective external fact. Caught by reference checks against authoritative sources. \textit{See Chapter 15.}
  \item[Fail-closed / fail-open.] Two designs for what an agent does when a check is uncertain. Fail-closed refuses the action; fail-open lets it through. For agents that move money, fail-closed is the only safe default. \textit{See Chapter 3.}
  \item[Faithfulness.] An evaluation criterion: does the agent's answer reflect what the source documents actually say, without distortion? A faithful answer can still be wrong (if the source is wrong); but a wrong-but-faithful answer is fixable, while an unfaithful one is hard to debug. \textit{See Chapter 15.}
  \item[Fine-tuning.] Taking a pre-trained model and continuing training on a narrower dataset to specialize it for a domain or behavior. Cheaper than training from scratch, riskier than prompting. \textit{See Chapter 2.}
  \item[Functional evaluation.] Measuring whether the agent completes the real-world task it was assigned, broken down by the legs of the agent loop: task outcome, tool-call correctness, multi-turn coherence, context retention, error recovery. \textit{See Chapter 15.}
  \item[GAIA.] A 2023 academic benchmark (Mialon et al.) testing whether an agent can complete tasks that are easy for humans and hard for AI — finding a specific file in a filesystem, following multi-step instructions across applications. Useful as a hard-mode sanity check; not enough on its own. \textit{See Chapter 10.}
  \item[Golden dataset.] A curated, versioned collection of test inputs paired with ground-truth outputs, with metadata rich enough to support diagnosis when something fails. The single most consequential investment a team can make in evaluation; the single most common thing they cut to save schedule. \textit{See Chapter 14.}
  \item[Governance triangle.] The three roles that hold certification together: the agent's owner, the certification review board, and the regulator (or its proxy). Each watches the other two. \textit{See Chapter 19.}
  \item[Graceful degradation.] When an agent under load or facing a tool outage falls back to a smaller, safer mode rather than failing hard. The agent's version of the airplane's "memorized standard turn." \textit{See Chapter 8; Chapter 15.}
  \item[Ground truth.] For a given input, the correct or validated output, established through expert review or authoritative source. The thing you are grading against. \textit{See Chapter 14.}
  \item[Groundedness.] An evaluation criterion: every claim in the agent's output can be traced back to a specific source document. Stronger than faithfulness — groundedness is yes/no per claim, not per output. \textit{See Chapter 15; Shuster et al. 2021.}
  \item[Guardrails.] The pre-2024 generic term for "things that stop the model from doing the wrong thing." The book splits the term into three more specific ideas — in-context guardrails, output filters, and tool gates — and prefers those. \textit{See Chapter 3.}
  \item[Hallucination.] A confidently stated, fluent, wrong answer from an LLM. The word is informal; production evaluation uses \textit{faithfulness}, \textit{groundedness}, \textit{factual hallucination}, and \textit{contextual hallucination} instead. \textit{See Chapter 1; Chapter 15.}
  \item[Hallucination rate.] The fraction of agent outputs that contain at least one hallucinated claim, measured against a golden dataset. Report it split by factual vs. contextual; the aggregate hides the diagnosis. \textit{See Chapter 15.}
  \item[Hardening.] Pillar 2 of ARIA. The practice of designing, constraining, and observing an agent so the consequences of its mistakes are bounded by what your organization can afford to lose. \textit{See Chapter 5.}
  \item[HELM (Holistic Evaluation of Language Models).] A multi-dimensional benchmark suite (Liang et al. 2022; Stanford CRFM, 2022 onward) that scores models on accuracy, calibration, robustness, fairness, bias, toxicity, efficiency, and other axes. The most comprehensive single source for cross-model comparison; not agent-specific. \textit{See Chapter 10.}
  \item[HITL (human-in-the-loop).] A design pattern where a human approves, reviews, or corrects an agent action before or after it is committed. The original alignment mechanism and still the most reliable. \textit{See Chapter 4.}
  \item[In-context guardrail.] A rule the agent reads at the start of every turn — typically embedded in the system prompt or fetched from a policy store. The runtime version of "do not refund more than \$1,000." \textit{See Chapter 3.}
  \item[Indirect prompt injection.] An attack in which the malicious instruction reaches the agent through a document it retrieves, a tool result it parses, or a web page it visits — not from the user. Described in detail by Greshake et al. 2023. The harder one to defend against because it does not come from the user at all. \textit{See Chapter 6.}
  \item[Inner alignment.] Does the trained model actually pursue the specified goal, in the situations it will actually face? The harder of the two alignment problems, because the answer is not always visible from the training data. \textit{See Chapter 1.}
  \item[ISO/IEC 42001:2023.] The international standard for an AI management system — a vendor-neutral framework an enterprise can adopt to govern its AI portfolio. Voluntary, but increasingly expected by procurement teams and useful as the spine of a certification program. \textit{See Chapter 17; Chapter 19.}
  \item[Jailbreak resistance.] The agent's ability to refuse manipulated prompts designed to get it to ignore its safety rules. A benchmark and evaluation target. \textit{See Chapter 16.}
  \item[Judge calibration.] The process of checking that an \textit{LLM-as-judge} agrees with human graders on a held-out sample. Without calibration, judge drift is invisible and corrosive. \textit{See Chapter 14.}
  \item[Kill switch.] A documented, tested way to disable an agent in production within seconds. The single most important hardening control no team builds until they have needed it once. \textit{See Chapter 3; Chapter 8.}
  \item[KYA (Know Your Agent).] An emerging discipline (analogous to KYC in financial services) by which an agent's owner is accountable for, at any moment, what data the agent sees, what tools it can call, what evaluations it has passed, and who approved its current deployment. Not yet a formal standard. \textit{See Chapter 17 §6.}
  \item[Latency (p50, p95, p99).] The 50th, 95th, and 99th percentile response times. The first number is what your team will quote; the last is what your users will feel. \textit{See Chapter 11.}
  \item[Leaderboard.] A public ranking of models or agents against a shared benchmark. The thing every CTO has tabs of open; almost always misleading without the context of which benchmark, what version, what conditions. \textit{See Chapter 9; Chapter 12.}
  \item[Leaderboard gaming.] Training, prompting, or scoring choices that improve leaderboard position without improving the agent's underlying capability — overfitting, contamination, or methodology arbitrage. The reason your team should not buy on leaderboard score alone. \textit{See Chapter 12.}
  \item[Least privilege (for agents).] Grant the agent only the smallest set of tools, scopes, and data it needs to do its job. The rule borrowed from network security; the rule most often broken by speed-to-launch agent teams. \textit{See Chapter 7.}
  \item[LLM (large language model).] A model — trained on huge amounts of text — that takes a sequence of words and produces the next one, then the next. The engine inside every agent in this book. \textit{See Chapter 1.}
  \item[LLM-as-judge.] Using a language model to grade the outputs of another model against a written rubric. Cheap, scalable, and worth exactly what you put into calibrating it. \textit{See Chapter 14; Zheng et al. 2023.}
  \item[Load testing.] Measuring an agent's behavior at expected and elevated traffic levels, looking for degradation curves. Distinct from stress testing, which deliberately overloads to find failure modes. \textit{See Chapter 11.}
  \item[MAESTRO.] The Cloud Security Alliance's \textit{Multi-Agent Environment, Security, Threat, Risk, and Outcome} methodology (CSA, 2025) for threat-modeling agentic systems. An alternative vocabulary to MITRE ATLAS and OWASP LLM Top 10 that your vendors and consultants may surface in 2025–26 decks. ATLAS and OWASP remain the more widely used baselines. \textit{See Chapter 6.}
  \item[MCP (Model Context Protocol).] Anthropic's emerging open standard (2024–25) for how agents call external tools. Think of it as USB-C for AI: it standardizes the handshake so any agent can talk to any tool server without custom glue. The hardening implication is that each MCP server becomes its own trust boundary. \textit{See Chapter 7.}
  \item[MITRE ATLAS.] An adversarial-tactics catalog for ML systems (MITRE Corporation, 2023 onward) — the ML version of ATT\&CK. Gives your security team a shared vocabulary for the attacks named in Chapter 6, expressible as technique IDs rather than marketing phrases. \textit{See Chapter 6.}
  \item[Model card.] A short, structured document published with a model, describing what it was trained on, what it is good at, what it is bad at, and what has been tested. Originally proposed by Mitchell et al. 2019; now a near-universal expectation for any model going into production. \textit{See Chapter 2; Chapter 18.}
  \item[Model Context Protocol.] See \textit{MCP}.
  \item[MTTD (mean time to detection).] The gap between an agent misbehaving and someone noticing. A high MTTD is the operational signature of an under-instrumented estate. \textit{See Chapter 8.}
  \item[MTTR (mean time to recovery).] The gap between noticing a misbehavior and being back in correct operation. Companion KPI to MTTD. \textit{See Chapter 8.}
  \item[MRR (mean reciprocal rank).] A retrieval-quality metric (Voorhees 1998): the average of 1/rank of the first relevant document across a set of queries. A score near 1.0 means the right document is almost always first. \textit{See Chapter 15.}
  \item[Multi-agent collusion.] When two or more agents in a system reinforce each other's mistakes or jointly arrive at a worse outcome than either would alone. Easier to set up than most teams expect; harder to debug. \textit{See Chapter 6.}
  \item[nDCG (normalized discounted cumulative gain).] A retrieval-quality metric (Järvelin \& Kekäläinen 2002) that rewards a system for placing more relevant documents higher in the ranking, discounted by position and normalized to 0–1. Useful when multiple documents are relevant and rank order across the top-k matters. \textit{See Chapter 15.}
  \item[NIST AI RMF 1.0.] The US National Institute of Standards and Technology's AI Risk Management Framework (Tabassi 2023) — voluntary, principles-based, organized around Govern / Map / Measure / Manage. The closest the US has to a federal AI standard. \textit{See Chapter 17.}
  \item[NIST AI 600-1 (GenAI Profile).] The 2024 NIST companion document to the AI RMF, specific to generative AI. Names the risks (confabulation, dangerous content, data leakage, information integrity, environmental harm) and the controls. Useful as a checklist when scoping a certification program. \textit{See Chapter 16; Chapter 17.}
  \item[NIST SP 800-207 (Zero Trust Architecture).] The 2020 NIST publication that defines Zero Trust as the principle that no actor — human, script, or agent — is trusted by default; every action is authenticated, authorized, and logged. \textit{See Chapter 5; Chapter 7.}
  \item[Non-deterministic check.] A grader that involves judgment: a human reviewer, an LLM acting as a judge, a probabilistic similarity metric. Flexible, expensive, harder to audit. Required for anything where there is more than one valid answer. \textit{See Chapter 14.}
  \item[Notified body.] Under the EU AI Act, an independent third-party organization authorized to conduct conformity assessments for certain categories of high-risk AI systems. \textit{See Chapter 19.}
  \item[Observability (for agents).] The combination of traces, logs, and evaluations that lets you answer "what did this agent do, and why?" after the fact. Distinct from traditional APM; the unit of observation is the agent turn, not the HTTP request. \textit{See Chapter 8.}
  \item[Offline / online evaluation.] Two of the three modes in the evaluation taxonomy. \textit{Offline} runs the agent against a curated, versioned test set before any production traffic touches it — the quality gate. \textit{Online} measures the agent on a sample of live production traffic, where the inputs are real and the consequences are real. \textit{See Chapter 13.}
  \item[Outcome-normalized cost.] Cost per achieved business outcome (a resolved ticket, a closed case, a correct decision), not cost per token or per call. A \$0.30-per-call agent that resolves 95\% of tickets first time is cheaper than a \$0.05-per-call agent that needs five tries. \textit{See Chapter 9; Chapter 11.}
  \item[Outer alignment.] Did we specify the goal correctly? The easier of the two alignment problems to define; routinely the harder one in practice, because "what we wanted" turns out to be plural and contradictory. \textit{See Chapter 1.}
  \item[OWASP LLM Top 10.] A community-maintained list (OWASP, 2023; updated 2025) of the top ten security risks specific to LLM-based applications: prompt injection, sensitive information disclosure, supply chain, data poisoning, improper output handling, excessive agency, system prompt leakage, vector and embedding weaknesses, misinformation, unbounded consumption. The threat catalog Chapter 6 maps to layers. \textit{See Chapter 6.}
  \item[\textit{Pass@k}.] The probability that an agent gets a task right in at least one of \textit{k} attempts. A useful "ceiling" metric; an optimistic "reliability" metric. \textit{See Chapter 10.}
  \item[\textit{Pass\textasciicircum{}k}.] The probability that an agent gets the same task right on every one of \textit{k} independent attempts. The reliability metric; the one you want for enterprise deployments. Introduced by Yao et al. 2024 in τ-bench. \textit{See Chapter 10; Chapter 15.}
  \item[PDP / PEP (Policy Decision Point / Policy Enforcement Point).] The two halves of a Zero Trust authorization stack. The PDP decides whether an action is allowed based on identity, scope, policy, and runtime signals; the PEP enforces the decision. In an agent system, the tool gateway is the PEP. \textit{See Chapter 7.}
  \item[Pre-cert binder.] The collection of artifacts (model card, eval reports, threat model, behavior contract, owner sign-off, SBOM, risk assessment, data lineage) that an agent needs before its certification process can even begin. Mostly things that should have existed already. \textit{See Chapter 18.}
  \item[Pre-training.] The expensive first phase of training, in which a model learns general language from very large text corpora. The phase Maya's vendor does, not the phase her enterprise does. \textit{See Chapter 2.}
  \item[Preference data.] Pairs of model responses, ranked by human raters as "A is better than B." The substrate RLHF and its successors train on. \textit{See Chapter 2.}
  \item[Prompt.] The text instructions given to an LLM, including system instructions, conversation history, and the current user input. The handle by which everything an agent does is set up. \textit{See Chapter 1.}
  \item[Prompt injection (direct).] An attack in which the user types something into the agent designed to override its system prompt — "ignore previous instructions and instead…" The 101-level agent attack. \textit{See Chapter 6.}
  \item[Prompt injection (indirect).] See \textit{indirect prompt injection}.
  \item[Propose → Decide → Enforce → Record.] The mental model for every agent tool call: the model proposes an action, the PDP decides whether to permit, the PEP enforces the decision, and an immutable audit entry records the proposal, decision, enforcement, and outcome. \textit{See Chapter 7.}
  \item[Quality evaluation.] The class of evaluation concerned with whether the agent's output was any good — faithful, well-structured, in scope, tonally appropriate — as distinct from whether it was functionally correct. \textit{See Chapter 13; Chapter 15.}
  \item[RAG (retrieval-augmented generation).] A pattern (Lewis et al. 2020) where the LLM is given relevant retrieved documents alongside the user query, so it can ground its answer in them. The pattern under almost every enterprise AI assistant currently in production. \textit{See Chapter 1 (foundation); Chapter 6; Chapter 7.}
  \item[Recertification trigger.] A documented condition (time elapsed, model upgraded, threshold breached, incident occurred) that forces an agent through certification again. Without triggers, certification is a one-time event; with them, it is a discipline. \textit{See Chapter 16; Chapter 20.}
  \item[Red-teaming.] The discipline of running deliberate adversarial attacks against an agent to surface failures that the eval suite misses. Manual, automated, or hybrid; standard practice in serious enterprise programs. \textit{See Chapter 16; Chapter 18; Perez et al. 2022.}
  \item[Reflexion / Self-Refine.] Two 2023 papers (Shinn et al.; Madaan et al.) on multi-pass refinement loops where the agent drafts, critiques, and revises its own output. Real gains, real costs (extra inference calls), and a tendency to amplify the model's biases without external grading. \textit{See Chapter 3.}
  \item[Refusal quality.] Whether the agent says "I can't do that" for the right reasons and not the wrong ones. Evaluated against a curated suite of safe-to-refuse and should-have-answered prompts. \textit{See Chapter 16.}
  \item[Regression benchmark.] A benchmark run on every release to catch a worse model before it ships. The agentic equivalent of a unit-test gate. \textit{See Chapter 12.}
  \item[Retrieval-augmented generation.] See \textit{RAG}.
  \item[Retrieval-grounded evaluation.] Measuring the retrieval step itself, separately from the generation step — checking whether the documents the agent received contained the information needed to answer correctly. Distinct from \textit{groundedness}, which measures the model's fidelity to whatever was retrieved. Standard metrics: \textit{MRR} and \textit{nDCG}. \textit{See Chapter 15.}
  \item[Reward hacking.] The agent finds a shortcut in the reward signal itself, exploiting how the score is computed rather than the thing the score was meant to measure. A common training-time alignment failure. \textit{See Chapter 1.}
  \item[Reward model.] A separately trained model that scores how good an agent's responses are, used to drive the RLHF loop. The reward model is itself a thing that can be miscalibrated. \textit{See Chapter 2.}
  \item[Risk-based tier.] A classification (typically low, medium, or high) that determines how heavy a certification process an agent needs. Borrowed from financial regulation; in agents, the EU AI Act sets the floor for some tiers via Annex III. \textit{See Chapter 17.}
  \item[RLHF (reinforcement learning from human feedback).] A training method (Christiano et al. 2017; Ouyang et al. 2022) in which human raters compare model outputs, and the model is fine-tuned to prefer the rated-better ones. The technique that made ChatGPT useful. \textit{See Chapter 2.}
  \item[Sandboxing.] Running a high-risk tool, agent, or workflow in an isolated execution environment with explicit network egress rules, filesystem scope, and resource budgets — so the blast radius of a misbehavior is bounded by the sandbox, not by the host system. The operational version of least privilege for agents that execute code or operate desktops. \textit{See Chapter 7.}
  \item[SBOM (software bill of materials).] A formal list of every component an agent ships with — model weights, datasets, libraries, prompts, MCP servers, tool integrations, retrieval indices. The artifact an audit team will ask for first. \textit{See Chapter 18.}
  \item[Secrets management (for agents).] How an agent gets the credentials it needs to call tools, with the smallest possible lifetime and scope. "The agent has the key" is the wrong default; rotated, scoped, audited tokens are the right one. \textit{See Chapter 7.}
  \item[Self-reflection.] Asking the model to evaluate its own response before sending it: \textit{"Here is your draft answer. Find three things wrong with it. Now revise."} Catches some avoidable errors at the cost of an extra inference call. \textit{See Chapter 3.}
  \item[Sleeper agent.] A model trained (in research, deliberately) to behave normally during evaluation but to act badly when a specific trigger appears in the prompt. Demonstrated by Hubinger et al. 2024. The argument that behavioral verification alone is not sufficient. \textit{See Chapter 4.}
  \item[Specification gaming.] The agent does exactly what you asked, and gets a result you did not want. "Minimize average wait time" becomes "hang up faster." \textit{See Chapter 1.}
  \item[SR 11-7.] US Federal Reserve supervisory guidance on \textit{model risk management}, issued 2011, that became the default framework for model governance inside US banks. SR 11-7 predates LLMs and AI agents and assumes a model with a fixed input and a fixed output; for an agent that has tools, memory, and the ability to be re-prompted, banks now layer ISO/IEC 42001 or NIST AI RMF controls on top of the SR 11-7 baseline. \textit{See Chapter 17 (Northwind Bank vignette).}
  \item[State drift exploitation.] An attack class in which the adversary nudges the agent's internal state — its memory, retrieved-context window, or accumulated tool outputs — across many turns until the agent reaches a configuration the attacker can exploit. Slower than direct prompt injection and harder to spot; the defense is bounded memory, periodic state resets, and anomaly detection on session-level patterns. \textit{See Chapter 6.}
  \item[Stress testing.] Pushing the agent beyond its expected load to find the failure mode, not the breaking point. Where load testing measures, stress testing breaks. \textit{See Chapter 11; Chapter 15.}
  \item[Success criterion.] A single, falsifiable statement about what the agent must do, attached to a measurable threshold. "Routes 95\% of incoming claims to the correct triage queue" is one; "demonstrates strong reasoning" is not. \textit{See Chapter 14.}
  \item[Supervised fine-tuning (SFT).] The training phase between pre-training and RLHF, in which a model is shown labeled examples of the kind of output you want it to produce. Cheaper than RLHF; less powerful at shaping preferences. \textit{See Chapter 2.}
  \item[SWE-bench.] A benchmark (Jimenez et al. 2024) where an agent tries to resolve real GitHub issues in twelve open-source Python repositories. The most-cited code-agent benchmark; widely misquoted by vendors who quote the original-version score, not the verified-version score. \textit{See Chapter 10.}
  \item[SWE-bench Verified.] A cleaned-up 500-instance subset of SWE-bench, released in August 2024 by OpenAI's team with the original authors, where every test was confirmed solvable and well-specified. The version any serious reference now means by "SWE-bench." \textit{See Chapter 10.}
  \item[Sycophancy.] When the agent tells the user what they want to hear, regardless of correctness. A common training-time alignment failure; documented across major models. \textit{See Chapter 1; Chapter 2.}
  \item[System card.] A model card's bigger cousin: a structured public document describing not just a model but an entire deployed system, including the tools it can call, the populations it serves, and the safeguards it operates under. \textit{See Chapter 18; Chapter 20.}
  \item[System prompt.] The instructions an agent reads at the start of every conversation, separate from anything the user types. The agent's job description and the first place runtime alignment lives. \textit{See Chapter 3.}
  \item[τ-bench.] Sierra's benchmark (Yao et al. 2024) for multi-turn conversational agents that have to elicit goals from a user. The benchmark that introduced \textit{pass\textasciicircum{}k}. The closest public benchmark to most enterprise customer-service use cases. \textit{See Chapter 10.}
  \item[Tabletop exercise.] A walkthrough drill in which a cross-functional team talks through a realistic agent incident — detection, containment, forensics, recovery, post-incident — without touching production. Borrowed from security and incident-response practice. The teams who run a tabletop before launch find the gaps before the auditor does. \textit{See Chapter 6; Chapter 8.}
  \item[Tail latency.] The slow end of a response-time distribution — typically the p95 and p99 percentiles. Most users will not register the difference between a 600 ms response and an 800 ms one, but they will notice a 22-second tail-event response. Tail latency is a product decision, not an engineering preference. \textit{See Chapter 11.}
  \item[Third-party assessor.] An independent assurance provider who reviews evidence and issues a certification decision. Required by some regulations (parts of the EU AI Act); optional under others (ISO/IEC 42001). \textit{See Chapter 19.}
  \item[Threshold / context / confidence gates.] Three descriptive sub-types of \textit{tool gate} used in Chapter 3. \textit{Threshold gates} permit an action when a numeric input (dollar amount, record count, blast-radius score) sits below a configured limit. \textit{Context gates} permit an action only when contextual signals — time of day, user role, source-of-request — are within an allowed set. \textit{Confidence gates} permit an action only when the model's self-reported confidence exceeds a calibrated floor; useful in narrow contexts and dangerous when the model is calibrated badly. \textit{See Chapter 3.}
  \item[Time-to-recover.] The duration from when an agent goes wrong (or alerts) to when it is back in correct operation. A KPI that becomes important after the first production incident, never before. \textit{See Chapter 15.}
  \item[Tool call.] The atomic act of an agent invoking a function, API, or service to do something in the world — search a database, send an email, refund a charge. The line between advisory intelligence and operational autonomy. \textit{See Chapter 1.}
  \item[Tool gate.] A runtime check between the agent's "I want to call tool X with arguments Y" and the tool actually executing. Where most of the operational guardrails for agents live. \textit{See Chapter 3.}
  \item[Tool gateway.] The architectural component that mediates every tool call: enforces scope, applies policy, logs the call, and returns the result. The PEP for agent tooling. \textit{See Chapter 7.}
  \item[Trace (agent trace).] A structured record of one agent turn — the prompt, the tools called, the arguments, the responses, the final output. The unit of observability and the unit of incident review. \textit{See Chapter 4; Chapter 8.}
  \item[Transparency.] What the agent owes its operator on every turn: a usable record of what it did and why. Distinct from explainability, which is what the agent owes a regulator after the fact. \textit{See Chapter 4; Chapter 20.}
  \item[Transparency report.] A recurring public summary of what the agent did and how it performed — what it refused, what it escalated, what failed. The reader is your customer, your regulator, or your civil-society watcher. \textit{See Chapter 20.}
  \item[Trust boundary.] The conceptual edge between two parts of a system that don't, by default, trust each other. In agent systems, every tool, every MCP server, and every retrieved document sits on the other side of a trust boundary from the agent. \textit{See Chapter 5; Chapter 7.}
  \item[TTFT (time to first token).] For streaming agents, the time from request to the first character of response. Distinct from total latency. A 6-second response that starts streaming at 400 ms feels fast; a 2-second response that stalls and then dumps the whole answer feels slow. \textit{See Chapter 11.}
  \item[Version-aware benchmarking.] Tagging every benchmark run with the model version, prompt version, tool versions, dataset version, and date — so when a number changes, you know what could plausibly have caused it. \textit{See Chapter 12.}
  \item[Workload identity.] The named, scoped, rotatable, revocable identity that every production agent should have — distinct from a shared service principal and distinct from any human's credentials. \textit{See Chapter 7.}
  \item[Zero Trust.] A security approach (NIST SP 800-207) in which no actor — inside or outside the network — is trusted by default; every request is authenticated, authorized, and logged. Applied to agents, every tool call is a Zero Trust event. \textit{See Chapter 5; Chapter 7.}
\end{description}

\stoptwocol
